\scanpage{16}

% head
\mspage{\textcircled{13}}
c) Conditions équivalentes

% condition (i)
(i) \inserted{$2G=0$ i.e. $G\simeq V$, et} $q(u)\ne0$ pour \deleted{\unsure{$u\ne0$}} $u\in V^*=V-\{0\}$

% condition (ii), ll. 3–6
(ii) Posant $\widetilde P_i=\widetilde E/\rho_i(G)$, \deleted{\unsure{la fon Pon|la fam Posons}} et $\widetilde p_i:\widetilde E\unsure{\longrightarrow|\mathbin{-}}\widetilde P_i$, si $x_1,x_2,x_3\in\widetilde\Gamma$ sont tels que $p_1(\widetilde x_2)=p_1(\widetilde x_3)$, $p_{\unsure{2|3}}(\widetilde z_3)=p_2(\widetilde x_{\unsure{1|3}})$, $p_3(\widetilde x_1)=p_3(\widetilde x_2)$, alors $\widetilde x_{\unsure{1|2}}=\widetilde x_2=\widetilde x_3$.

% condition (iii), ll. 7–8
(iii) Si $\widetilde z_1\in\rho_1(G)$, $\widetilde z_2\in\rho_2(G))$, $\widetilde z_3\in\rho_3(G)$ sont tels que $\widetilde z_1\widetilde z_2\widetilde z_3=0$, alors $\widetilde z_1=\widetilde z_2=\widetilde z_3=1$

% condition (iv), ll. 9–10
(iv) Si $g\in G$ est tel que $\rho_1(g)\rho_2(g)\rho_3(g)=1$, alors $g\unsure{=|\equiv}1$

% conclusion, ll. 11–14
De plus, pour $G,\varepsilon$ fixés avec $\varepsilon\unsure{\ne|=}\mathbb F_2$, $\exists$ une $\widetilde\Gamma$ \deleted{ext} avec ces propriétés \unsure{sss|ssi} $2G=0$ i.e. $G\simeq V$, $\operatorname{rg}_{\mathbb F_2}V\leq2$ (et cette extension est unique à isom unique près).
% archive sheet number: 15
